% Standalone problem document, generated from the adjacent README.md.
% Compile with XeLaTeX. Mathematical content is maintained in Markdown.
\documentclass[11pt,a4paper]{article}
\usepackage[fontsize=10.5pt]{scrextend}
\usepackage[margin=22mm,headheight=15pt,headsep=9mm,footskip=11mm]{geometry}
\usepackage{fontspec}
\setmainfont{texgyrepagella}[Extension=.otf,UprightFont=*-regular,BoldFont=*-bold,ItalicFont=*-italic,BoldItalicFont=*-bolditalic]
\setsansfont{texgyreheros}[Extension=.otf,UprightFont=*-regular,BoldFont=*-bold,ItalicFont=*-italic,BoldItalicFont=*-bolditalic]
\setmonofont{texgyrecursor}[Extension=.otf,UprightFont=*-regular,BoldFont=*-bold,ItalicFont=*-italic,BoldItalicFont=*-bolditalic,Scale=MatchLowercase]
\usepackage{amsmath,amssymb}
\usepackage{unicode-math}
\setmathfont{texgyrepagella-math.otf}
\usepackage{xcolor,microtype,enumitem,needspace,fancyhdr}
\usepackage{bookmark}
\definecolor{ink}{HTML}{17344B}
\definecolor{accent}{HTML}{197D80}
\definecolor{muted}{HTML}{596773}
\hypersetup{colorlinks=true,linkcolor=accent,urlcolor=accent,pdftitle={MI-27 - Constant
one in the logarithmic commutator
inequality},pdfauthor={Open Problems in Numerical Linear Algebra}}
\urlstyle{same}
\setlength{\parindent}{0pt}
\setlength{\parskip}{5pt plus 1pt minus 1pt}
\linespread{1.04}
\setlength{\emergencystretch}{3em}
\widowpenalty=10000
\clubpenalty=10000
\displaywidowpenalty=10000
\allowdisplaybreaks[1]
\setlist{nosep,leftmargin=1.5em}
\providecommand{\tightlist}{\setlength{\itemsep}{0pt}\setlength{\parskip}{0pt}}
\providecommand{\pandocbounded}[1]{#1}
\pagestyle{fancy}
\fancyhf{}
\fancyhead[L]{\sffamily\footnotesize\color{muted}OPEN PROBLEMS IN NUMERICAL LINEAR ALGEBRA}
\fancyhead[R]{\sffamily\bfseries\footnotesize\color{accent}MI-27}
\fancyfoot[L]{\parbox[b]{0.9\textwidth}{\sffamily\fontsize{8}{10}\selectfont\color{muted}Editorial ratings; a bounded search cannot certify that a problem remains unsolved.\\Literature check: 2026-09-12}}
\fancyfoot[R]{\sffamily\footnotesize\color{muted}\thepage}
\renewcommand{\headrulewidth}{0.3pt}
\renewcommand{\headrule}{\hbox to\headwidth{\color{accent}\leaders\hrule height \headrulewidth\hfill}}
\makeatletter
\renewcommand{\section}{\Needspace{5\baselineskip}\@startsection{section}{1}{0pt}{12pt plus 2pt minus 2pt}{4pt}{\sffamily\bfseries\large\color{ink}}}
\renewcommand{\subsection}{\Needspace{4\baselineskip}\@startsection{subsection}{2}{0pt}{9pt plus 1pt minus 1pt}{3pt}{\sffamily\bfseries\normalsize\color{ink}}}
\makeatother
\setcounter{secnumdepth}{0}
\begin{document}
{\sffamily\small\color{muted}Matrix inequalities and norms\par}
\vspace{3pt}
{\raggedright\hyphenpenalty=10000\exhyphenpenalty=10000\sffamily\bfseries\fontsize{21}{24}\selectfont\color{ink}Constant
one in the logarithmic commutator inequality\par}
\vspace{7pt}
{\sffamily\small\textbf{Difficulty:} challenging\quad\textbf{Importance:} broadly
interesting\par}
{\sffamily\small\bfseries\color{accent}Status: Solved\par}
\vspace{4pt}
{\color{accent}\hrule height 0.8pt}
\vspace{6pt}
\textbf{Historical rating rationale:} Improving the known universal
coefficient to its proposed sharp value is challenging; connections
between matrix logarithms, entropy and quantum dynamics make the
question broadly interesting.

\subsection{Resolution: affirmative, 12 September
2026}\label{resolution-affirmative-12-september-2026}

\textbf{Sidney Holden}, Center for Computational Biology, Flatiron
Institute, Simons Foundation, proves the exact inequality for every
dimension and every complex positive definite pair in the original
statement, with optimality already in order two. See \textbf{Theorem
1.1, Sections 2--5}, of the
\href{https://github.com/ajt60gaibb/OpenProblemsInNLA/blob/main/references/holden-mi27-2026-09-12/solution.pdf}{complete
manuscript}
(\href{https://github.com/ajt60gaibb/OpenProblemsInNLA/blob/main/references/holden-mi27-2026-09-12/solution.md}{Markdown},
\href{https://github.com/ajt60gaibb/OpenProblemsInNLA/blob/main/references/holden-mi27-2026-09-12/solution.tex}{LaTeX}).

A
\href{https://github.com/ajt60gaibb/OpenProblemsInNLA/blob/main/references/holden-mi27-2026-09-12/verification/independent-review.md}{separate
Codex AI-agent audit} passed the full proof and checked its published
relative-entropy input. This is informal review; no external human peer
review, Lean verification or priority claim is asserted.
\href{https://github.com/ajt60gaibb/OpenProblemsInNLA/blob/main/references/holden-mi27-2026-09-12/provenance.md}{Provenance
and verified affiliation} disclose AI assistance and credit
\href{https://github.com/ajt60gaibb/OpenProblemsInNLA/pull/186}{PR
\#186}, whose partial results left the universal upper bound open. The
original target below is retained; ratings above are historical.

\subsection{Problem statement}\label{problem-statement}

For every integer \(n\ge1\) and positive definite
\(A,B\in\mathbb C^{n\times n}\) with
\(\mathop{\mathrm{tr}}\nolimits(A+B)=1\), set
\(a=\mathop{\mathrm{tr}}\nolimits A\),
\(b=\mathop{\mathrm{tr}}\nolimits B\). Is

\[\bigl\|B\log(A+B)-\log(A+B)B\bigr\|_1
\le -a\log a-b\log b\]

always true? The logarithm is the natural logarithm applied by spectral
functional calculus;
\(\|X\|_1=\mathop{\mathrm{tr}}\nolimits(X^*X)^{1/2}\). Thus \(a,b>0\)
and \(a+b=1\), and every expression is well-defined.

The target is specifically the coefficient \(1\) on the right. Existence
of some dimension-independent coefficient has been proved and is not the
open question. The coefficient-one question is explicitly suggested
immediately after the original displayed conjecture.

This is a sharp matrix-function estimate for the rate at which
noncommuting positive matrices mix. It connects perturbation estimates
for the matrix logarithm with entropy production and quantum
information.

\newpage
\subsection{References}\label{references}

\begin{enumerate}
\def\labelenumi{\arabic{enumi}.}
\item
  K. M. R. Audenaert and F. Kittaneh, \emph{Problems and Conjectures in
  Matrix and Operator Inequalities}, arXiv:1201.5232v3 (2012), §6,
  Conjecture 4, equation (26), and the sentence proposing \(c=1\).
  \href{https://arxiv.org/html/1201.5232}{Full text}.
\item
  K. M. R. Audenaert, \emph{Quantum Skew Divergence}, J. Math. Phys. 55
  (2014), 112202, §8, proof of small incremental mixing with coefficient
  \(2\). \href{https://arxiv.org/html/1304.5935}{Full text}.
\item
  A. Vershynina, \emph{Entanglement rates for Renyi, Tsallis and other
  entropies}, arXiv:1803.07117v3 (2021), §8, conclusion, pp.~15--16,
  explicitly distinguishing the proved coefficient \(2\) and proposed
  coefficient \(1\). \href{https://arxiv.org/pdf/1803.07117}{Full
  preprint}.
\item
  Q. Ning, F.-Z. Guo, J. Zhang and Q.-Y. Wen, \emph{On Bounding
  Entangling Rates and Mixing Rates in Some Special Cases}, Int. J.
  Theor. Phys. 55 (2016), 1686--1694.
  \href{https://doi.org/10.1007/s10773-015-2806-9}{Published abstract},
  which reports coefficient one under additional restrictions; the
  subscription body was not independently inspected in this audit.
\end{enumerate}

Historical status check (2026-09-10; superseded by the resolution
above): Audenaert's theorem settles the existence claim with \(2\).
Vershynina's 2021 revision distinguishes the proved coefficient from the
proposed coefficient \(1\). Searches for ``small incremental mixing
sharp constant'', ``logarithmic commutator Audenaert constant one'', and
2025/2026 found later special-case estimates but no announced universal
coefficient-one proof or counterexample. This limited search does not
certify that the conjecture remains open.
\end{document}
